\section{Descubrimiento de joins en el agente de validación}
\label{sec:join-discovery}

Muchos conjuntos de datos reales no llegan como un único fichero, sino repartidos en
varias tablas relacionadas: un esquema en estrella con su tabla de hechos y sus
dimensiones, una serie temporal acompañada de tablas de enriquecimiento (calendario,
meteorología, catálogo de productos), etc. El \textit{pipeline}, sin embargo, opera
sobre un \textbf{dataset canónico} único. Cerrar esa brecha es responsabilidad del
agente de validación de datos (\texttt{data\_validator}), que cuando recibe varios
ficheros ejecuta un subproceso de \textbf{descubrimiento de \textit{joins} agéntico}:
infiere qué columnas enlazan las tablas y las integra en una sola, \emph{sin} que el
usuario tenga que declarar las claves de \textit{join}.

Esta sección detalla el mecanismo completo. La figura~\ref{fig:join-discovery-flow}
ofrece la panorámica del subproceso dentro del nodo; aquí se desarrollan los pasos, las
métricas deterministas y las garantías de cada herramienta.

\subsection{Principio de diseño: separar la decisión de la medición}
\label{sec:join-principio}

El descubrimiento de \textit{joins} encarna el principio de \textit{restraint} agéntico
que guía todo el sistema: el modelo de lenguaje interviene únicamente donde aporta
valor —reducir un espacio de búsqueda combinatorio y emitir un juicio semántico sobre
qué columnas son claves plausibles— mientras que toda medición y toda transformación de
datos las realiza código determinista.

Concretamente, el agente solo hace dos cosas: \textbf{(i)} proponer un conjunto reducido
de candidatos de \textit{join} razonando sobre perfiles estructurales, y \textbf{(ii)}
tomar la decisión final de aceptar o rechazar cada candidato. No calcula ninguna métrica
de calidad ni ejecuta ningún \textit{merge}. Como consecuencia, el agente \emph{no puede
alucinar} la calidad de un \textit{join}: cualquier selección debe estar respaldada por
las cifras que devuelve la herramienta determinista de evaluación, y la materialización
del \textit{join} es reproducible e independiente del modelo. Esta división detección
determinista / decisión agéntica es la misma que gobierna la validación de datos y la
planificación del entrenamiento.

\subsection{Perfilado determinista previo}
\label{sec:join-perfilado}

El subproceso se activa solo cuando se proporciona más de un fichero. Conviene subrayar
que el perfilado \textbf{no es una herramienta del agente}, sino código del nodo que se
ejecuta \emph{antes} de que el agente arranque: el agente no puede invocarlo y, por tanto,
no lo calcula. La función \texttt{profile\_raw\_datasets} —Python puro, sin decorador de
herramienta ni llamadas al modelo— se ejecuta una sola vez durante el ensamblaje del
mensaje de entrada (el \texttt{HumanMessage} que recibe el \texttt{data\_validator}), de
modo que los perfiles viajan \textbf{embebidos en ese propio mensaje de entrada}. Las
únicas herramientas de \textit{join} que el agente sí puede invocar son
\texttt{evaluate\_join\_candidates} y \texttt{execute\_join\_plan}. Para cada columna de cada tabla se calcula un perfil
compacto: tipo de dato, número de valores no nulos, tasa de nulos, cardinalidad
(\texttt{unique\_count}), ratio de unicidad (\texttt{unique\_ratio}) y, para columnas
numéricas o de fecha, los valores mínimo y máximo. Se adjuntan además las primeras
\texttt{data\_validator\_profile\_nrows} filas de muestra (10 por defecto).

Estos perfiles se serializan e \textbf{inyectan en el mensaje inicial} del agente. De
este modo, el agente dispone del contexto estructural de todos los ficheros —nombres de
columnas, tipos, cardinalidades, rangos y filas de ejemplo— sin haber invocado todavía
ninguna herramienta y sin que el modelo haya visto los datos completos. El perfilado
recorre los ficheros enteros para que las estadísticas sean exactas, pero solo expone
una muestra acotada, lo que mantiene el coste en \textit{tokens} bajo y constante con
independencia del tamaño de los ficheros.

\subsection{Propuesta de candidatos por el agente}
\label{sec:join-candidatos}

A partir de los perfiles inyectados, el agente razona en el bucle ReAct y produce una
propuesta en dos partes:

\begin{enumerate}
  \item \textbf{Selección del \textit{dataset} base.} El agente elige la tabla base, cuyo
        número de filas \emph{define la cardinalidad del dataset canónico final}: todas
        las demás tablas se enlazarán a ella mediante \textit{left join}, rellenando con
        \texttt{NA} las filas sin correspondencia. El criterio es la cobertura de las
        columnas objetivo; en caso de empate se prefiere la tabla con más filas.
  \item \textbf{Propuesta de candidatos.} El agente identifica qué columnas del esquema
        objetivo no están cubiertas por la tabla base y propone \textbf{como máximo 20}
        candidatos de \textit{join} —pares (tabla.columna izquierda, tabla.columna
        derecha)— que podrían aportarlas. La instrucción es explícita: debe razonar sobre
        los perfiles (nombres, tipos, \texttt{unique\_ratio}, rangos, muestras) para
        proponer claves \emph{semánticamente plausibles}, no adivinar a ciegas.
\end{enumerate}

El papel del agente en este paso es esencialmente \textbf{podar el espacio de búsqueda}:
en lugar de evaluar el producto cartesiano de todas las parejas de columnas posibles,
selecciona un puñado de hipótesis razonables que el código evaluará a continuación.

\subsection{Evaluación determinista de candidatos}
\label{sec:join-evaluacion}

El agente invoca entonces la herramienta \texttt{evaluate\_join\_candidates}, que para
cada candidato propuesto mide su calidad de forma determinista. Los valores de las claves
se normalizan (\textit{trim} y minúsculas) antes de comparar, para tolerar diferencias de
formato. La tabla~\ref{tab:join-metricas} resume las métricas calculadas y su papel en la
decisión.

\begin{table}[H]
\centering
\caption{Métricas deterministas calculadas por \texttt{evaluate\_join\_candidates} para
cada candidato de \textit{join}. Los umbrales son configurables; se indican los valores
por defecto.}
\label{tab:join-metricas}
\small
\setlength{\tabcolsep}{5pt}
\begin{tabular}{lp{6.2cm}l}
\toprule
\textbf{Métrica} & \textbf{Definición} & \textbf{Umbral} \\
\midrule
\texttt{left\_coverage}  & Fracción de claves distintas del lado izquierdo con
                           correspondencia en el derecho ($|L\cap R| / |L|$). & $\geq 0{,}8$ \\
\texttt{right\_coverage} & Fracción análoga del lado derecho ($|L\cap R| / |R|$). & --- \\
\texttt{containment}     & $\max$(cobertura izq., der.); detecta relación de subconjunto
                           (clave foránea $\subseteq$ clave primaria). & $\geq 0{,}8$ \\
\texttt{jaccard}         & Solapamiento simétrico $|L\cap R| / |L\cup R|$. & --- \\
\texttt{inferred\_rel.}  & Relación inferida (1:1, 1:N, N:1, N:N) según el
                           \texttt{unique\_ratio} de cada lado. & UR $\geq 0{,}98$ \\
\texttt{row\_multiplier} & Filas estimadas tras el \textit{join} dividido por las filas
                           base; estima la \emph{explosión de filas}. & --- \\
\texttt{explosion\_risk} & Riesgo derivado del multiplicador: bajo, medio o alto. &
                           $1{,}25$ / $2{,}0$ \\
\bottomrule
\end{tabular}
\end{table}

Además de las métricas, la herramienta emite \textbf{advertencias} explícitas cuando
detecta condiciones problemáticas: discrepancia de tipos entre las dos columnas clave,
cobertura izquierda o contención por debajo del umbral, relación \textit{many-to-many}, o
presencia de nulos en alguna de las claves. Para acotar el coste, rechaza de entrada
cualquier llamada con más de 20 candidatos. El resultado es una lista estructurada de
evaluaciones que el agente recibe y sobre la que fundamentará su decisión.

\subsection{Selección del plan y ejecución del \textit{join}}
\label{sec:join-ejecucion}

Con las evaluaciones en mano, el agente compone un \texttt{JoinPlan} estructurado:
\textit{joins} seleccionados (con las columnas a añadir, una confianza y una
justificación), candidatos rechazados con motivo explícito, ambigüedades sin resolver y
advertencias. El sistema le obliga a rechazar un candidato cuando ninguno alcanza
cobertura o contención aceptables, cuando la relación es \textit{many-to-many} o cuando el
riesgo de explosión de filas es alto. La regla de oro es que \emph{no puede seleccionar un
\textit{join} que la evaluación no haya confirmado}.

La materialización corre a cargo de \texttt{execute\_join\_plan}, también determinista,
que ofrece varias garantías sobre el dataset canónico resultante:

\begin{itemize}
  \item \textbf{Solo \textit{left join}.} El \textit{dataset} base fija las filas; su
        cardinalidad se preserva siempre. Las filas de enriquecimiento sin
        correspondencia se rellenan con \texttt{NA}. No se admiten \textit{inner joins}.
  \item \textbf{Deduplicación de la clave derecha.} Antes de cada \textit{merge} se
        eliminan duplicados de la clave en la tabla derecha (se retiene la primera
        aparición), evitando una explosión accidental de filas; cada deduplicación se
        registra como advertencia.
  \item \textbf{Resolución de colisiones de columnas.} Las columnas que ya existen en la
        tabla en curso se renombran con el sufijo del \textit{dataset} de origen, y la
        clave derecha redundante se descarta tras el \textit{merge}.
  \item \textbf{Bloqueo de claves erróneas.} Si un candidato seleccionado tiene
        solapamiento nulo, la ejecución se aborta con un error explícito (par de claves
        incorrecto), impidiendo producir un dataset vacío de forma silenciosa.
  \item \textbf{Verificación del esquema.} Tras todos los \textit{joins}, comprueba que
        las columnas objetivo requeridas estén presentes; si falta alguna, falla de forma
        controlada en lugar de continuar con un dataset incompleto.
\end{itemize}

El fichero combinado se escribe en disco y se trata como un único \textit{dataset}
canónico para el resto del \textit{pipeline}. La herramienta devuelve además el
\texttt{JoinPlan} completo, que el nodo escribe en el estado compartido
(\texttt{data\_join\_plan}).

\subsection{Auditoría y supervisión humana}
\label{sec:join-auditoria}

Como todo \textit{join} altera la semántica del \textit{dataset}, el plan no se aplica de
forma opaca. El \texttt{JoinPlan} auditado viaja en el estado hasta la primera puerta de
aprobación humana (\texttt{dataset\_approval}), donde el operador ve, junto al informe de
validación y la vista previa del dataset procesado, las claves de \textit{join}
empleadas, las métricas de cobertura y cardinalidad por candidato, y los candidatos
rechazados con su motivo. Así puede verificar no solo que el resultado es correcto, sino
que las decisiones de integración del agente son razonables. Si el operador rechaza con un
comentario, el controlador lo inyecta como retroalimentación en el estado y reencamina al
agente para que intente un enfoque alternativo.

El comportamiento de este subproceso sobre casos reales se examina en el capítulo de
resultados: el descubrimiento de tres \textit{joins} multi-tabla en \textit{Grid Demand}
(Sección~\ref{sec:caso-joins}) y el manejo de \textit{joins} imperfectos —cobertura
parcial, claves foráneas huérfanas, deduplicación— en \textit{retail\_sales}
(Sección~\ref{sec:caso-joins-imperfectos}).
